
A conversational agent that cannot remember is not a companion; it is a
sequence of strangers. Yet remembering is bounded by the context window: a model
can only attend to what fits, and a relationship that spans weeks does not fit.
The standard response is retrieval-augmented memory — extract what matters, store
it, retrieve it on demand — and a growing line of work (MemGPT, Mem0, Generative
Agents) has shown this can give agents durable, useful memory.

That work shares an assumption we want to remove: **the memory manager is a
large, capable model**, usually reached through an API. For a companion that
runs on a user's own machine — private, offline, always-available — this is
exactly the wrong dependency. The open question is not whether memory helps (it
does), but whether the *management* of memory — deciding what to store, how to
date it, when to update or forget — can be done by a small model on consumer
hardware, and if so, how small.

We study this concretely on a single 12 GB GPU, with the chat model held fixed
and only the memory pathway varied. We compare an \textbf{agentic} memory protocol —
which extracts typed facts (durable facts, dated events, the agent's own
self-facts), absolutizes relative dates at write time, deduplicates on insert,
and retrieves age-aware — against three baselines: no memory, stuffing the full
transcript into the window, and retrieving raw dialogue turns without extraction.

Our findings reframe how these strategies should be compared. Judged on
aggregate accuracy, full-history stuffing wins — but we show this win is
\textbf{contingent on conversation length}. Splitting questions by whether their
evidence still fits in the window, stuffing's accuracy collapses 76\% once the
evidence is truncated, while retrieval-based memory is \textbf{length-invariant},
and overtakes stuffing 2.0× in exactly the long-conversation regime a persistent
companion lives in. Separately, extracting and *dating* memories at write time
yields the best temporal reasoning of any strategy, at ~67× less context — the
result that makes local deployment practical. Finally, sweeping the extractor
from 3B to 14B, we find a \textbf{minimum viable size}: below ~7B the typed protocol
fails, and temporal grounding is the first thing lost.

\textbf{Contributions.}
1. A reproducible benchmark of an agentic FACT/EVENT/SELF/UPDATE/FORGET memory
   protocol against no-memory, full-history, and raw-retrieval baselines, judged
   and released with raw outputs.
2. The \textbf{length-invariance} result: context-stuffing's accuracy advantage
   expires as conversations grow, while retrieval-based memory does not — with a
   direct measurement of the crossover.
3. An answer to \textbf{how small the memory manager can be} (≈7–14B), identifying
   temporal grounding as the capability that breaks first under compression.
4. An open-source system running the full recall–generate–extract loop in real
   time on a single 12 GB consumer GPU.

---
